\section{Related Work}
\label{sec:related}

\paragraph{Semantic query processing.}
Semantic operators extend relational processing with natural-language filters,
joins, maps, rankings, and aggregations. LOTUS formalizes these operators and
uses proxy models, embeddings, and statistical procedures to reduce cost while
providing accuracy guarantees \cite{patel2025lotus}. Palimpzest searches plans
across models, prompts, and physical implementations
\cite{liu2025palimpzest}; UQE combines foundation-model calls with sampling and
compiler techniques \cite{dai2024uqe}; and DocETL uses agentic rewrites and
plan evaluation for complex document pipelines \cite{shankar2025docetl}.
SemBench provides a common evaluation of semantic query engines across
operators and modalities \cite{lao2026sembench}. \system focuses on a narrower
physical decision inside these systems---which answer model executes each
operator instance---and on the prompt/data statistics required to make that
decision reliable. Unlike operator-level plan search, the route is
instance-specific and rare failures dominate its risk.

\paragraph{LLM cascades and routers.}
FrugalGPT studies prompt adaptation, approximation, and cascades across model
APIs \cite{chen2024frugalgpt}. Language Model Cascades learns post-hoc deferral
from token-level uncertainty \cite{gupta2024cascades}. RouteLLM trains several
routers from preference data and introduces cost--performance metrics
\cite{ong2025routellm}; UniRoute represents models in a common feature space so
new models can be routed without retraining \cite{jitkrittum2026uniroute}.
These works establish the potential of learned routing, often on large
preference or benchmark collections. Our study instead examines low-label,
operator-specific routing where both answer models can be executed offline and
correctness labels are available. The four-outcome abstraction exposes two
cases hidden by strong-versus-weak preference labels: Nano can be uniquely
correct, and both models can fail. It also makes fresh critical recall a
first-class measure rather than inferring reliability from an aggregate curve.

\paragraph{Automatic and soft prompt optimization.}
ProTeGi constructs textual gradients and searches prompt edits
\cite{pryzant2023protegi}; MIPRO jointly optimizes instructions and
demonstrations \cite{opsahlong2024mipro}; DSPy compiles declarative LM programs
against downstream metrics \cite{khattab2024dspy}; and GEPA uses natural-language
reflection, execution traces, and Pareto-aware evolution
\cite{agrawal2025gepa}. Recent analysis shows that reflective optimization can
follow uninterpretable trajectories and fail systematically when its feedback
or evaluation signal is weak \cite{wang2026reflection}. Our results provide a
systems instance of this warning: richer routing prose does not prevent an
all-Nano collapse, and validation-calibrated scores miss fresh rare events.
Soft prompt tuning instead optimizes continuous prefix embeddings while
freezing most model parameters \cite{lester2021prompttuning}. We compare these
surfaces under one utility rather than treating automatic prompt engineering
and learned routing as separate problems.

\paragraph{Selective prediction and learning to defer.}
The constrained objective in Eq.~\ref{eq:constrained} is related to selective
classification and learning to defer, where a predictor abstains or transfers
an example to an expert \cite{madras2018defer}. In model routing, deferral has a
monetary cost and the expert can also be wrong. Our \no and \bw outcomes make
that distinction explicit, while the validation threshold supplies an
operating point for a query optimizer's quality target.
